\VerdierPhysicalPageAnchor{0058}%
\VerdierSemanticAnchor{ch1-1-5-10-converse}
\noindent\textbf{Conversely:}\par

\VerdierSemanticAnchor{ch1-1-5-10}
\VerdierSemanticAnchor{ch1-1-5-10-proposition}
\noindent\textbf{Proposition 1.5.10.}\quad
\emph{Every morphism $u\colon F_1\to F_2$ between the underlying
functors such that, for every $i\in[1,n]$, diagram (1.5.9.2) is
$\beta(e_i)$-commutative is a $\beta$-morphism of functors.}

\VerdierSemanticAnchor{ch1-1-5-10-proof}
Indeed, for every $\sigma\in\mathbb Z^n$, let $R(\sigma)$ denote the
property:

\VerdierSemanticAnchor{ch1-1-5-10-property}
\noindent\emph{$(R(\sigma))$\qquad The following diagram is
$\beta(\sigma)$-commutative:}

\VerdierSemanticAnchor{ch1-1-5-10-eq1}
\begin{equation}
  \resizebox{0.78\textwidth}{!}{%
    $\begin{tikzcd}[ampersand replacement=\&,column sep=6em,row sep=4em]
      F_1T(\sigma)
        \arrow[r,"{p_1(\sigma)}"]
        \arrow[d,"{u*T(\sigma)}"']
        \& T'(\sigma)F_1
          \arrow[d,"{T'(\sigma)*u}"] \\
      F_2T(\sigma)
        \arrow[r,"{p_2(\sigma)}"]
        \& T'(\sigma)F_2
    \end{tikzcd}$%
  }\;.
  \tag*{\ensuremath{(1.5.10.1)_{\sigma}}}
\end{equation}

Inspection of the various diagrams
$(1.5.10.1)_{\sigma_1}$, $(1.5.10.1)_{\sigma_2}$, and
$(1.5.10.1)_{\sigma_1+\sigma_2}$ immediately shows that:

\VerdierSemanticAnchor{ch1-1-5-10-closure}
\noindent\emph{If two of the properties $R(\sigma_1)$,
$R(\sigma_2)$, and $R(\sigma_1+\sigma_2)$ hold, then so does the
third.}

\VerdierSemanticAnchor{ch1-1-5-10-conclusion}
Moreover, by hypothesis, the properties $R(e_i)$, $i\in[1,n]$, and
$R(0)$ hold. It follows immediately that $R(\sigma)$ holds for every
$\sigma\in\mathbb Z^n$.

\VerdierSemanticAnchor{ch1-1-5-11}
\noindent\textbf{1.5.11.}\quad
A strict $\mathbb Z^n$-category is a category equipped with a
representation of $\mathbb Z^n$ in its group of automorphisms. To define
such a representation, it of course suffices to specify $n$
automorphisms:
$T(e_1),\ldots,T(e_n)$ which commute pairwise.

\VerdierSemanticAnchor{ch1-1-5-12}
\noindent\textbf{1.5.12.}\quad
The problem of determining on a given category the structures of a
$\mathbb Z^n$-category (not necessarily strict) is much less simple.
Thus, if we restrict ourselves to determining, on a category
$\mathcal C$, the $\mathbb Z^n$-category structures such that
$T(\sigma)=\text{the identity on }\mathcal C$ for every
$\sigma\in\mathbb Z^n$, we must then specify, for every pair
$\sigma_1,\sigma_2$ of elements of $\mathbb Z^n$, an element
$c(\sigma_1,\sigma_2)$ of the automorphism group of the identity
functor of $\mathcal C$ (this group is commutative), this datum being
required to satisfy identities (1.2.1.3). Consequently, to specify such
a structure, one must specify a normalized inhomogeneous 2-cocycle on
$\mathbb Z^n$ with values in the automorphism group of the identity
functor of $\mathcal C$. One also verifies immediately that two such
structures are isomorphic, by a $\mathbb Z^n$-functor $(F,m)$ such that
$F=\text{the identity on }\mathcal C$, if and only if the corresponding
cocycles lie in the same cohomology class.
